\section{Preliminaries}
\label{sec:preliminaries}

In this section, we recall background to understand the paper.
We give notation used throughout the paper as shown in Table~\ref{tab:notation}.

{\footnotesize
\setlength{\tabcolsep}{4pt}
\renewcommand{\arraystretch}{1.12}
\begin{longtable}{@{}>{\raggedright\arraybackslash}p{0.27\linewidth}
                  >{\raggedright\arraybackslash}p{0.67\linewidth}@{}}
\caption{Notation used throughout the paper.}
\label{tab:notation}\\
\toprule
\textbf{Notation} & \textbf{Meaning} \\
\midrule
\endfirsthead

\toprule
\textbf{Notation} & \textbf{Meaning} \\
\midrule
\endhead

\midrule
\multicolumn{2}{r@{}}{\emph{Continued on the next page}}\\
\endfoot

\bottomrule
\endlastfoot

\multicolumn{2}{@{}l}{\textbf{Basic parameters and finite fields}}\\
\midrule
\(\lambda\) & Security parameter. \\
\(p\) & Prime $p=c\cdot 2^f-1$ where $p\approx2^{2\lambda}$ and $c$ is the smallest cofactor. \\
\(\Fp, \Fpp\) & Finite fields with \(p\) and \(p^2\) elements, respectively. \\
\(K, \overline{K}\) & A finite field \(K\), and its algebraic closure. \\
% \(\Z,\Q,\N\) & The set of integers, rational numbers, and positive integers. \\
\((a:N)\) & Legendre symbol \(\left(\frac{a}{N}\right)\) modulo the odd prime \(N\). \\
% \(a\equiv b \pmod I\) & Congruence modulo the lattice \(I\). \\
% In algorithms, reduction modulo \(N\) is taken coefficient-wise in the fixed \(\calO_0\)-basis. \\

\addlinespace[3pt]
\multicolumn{2}{@{}l}{\textbf{Elliptic curves and isogenies}}\\
\midrule
% \(E,E',E''\) & Supersingular elliptic curves over \(\Fpp\). \\
\(E_0\) & Supersingular elliptic curve \(E_0:y^2=x^3+x\). \\
\(0_E\) & Identity point of the elliptic curve \(E\). \\
\([n]\) & Multiplication-by-\(n\) map on an elliptic curve. \\
% \(E[n]\) & \(n\)-torsion subgroup of \(E\). \\
% \(\varphi,\psi,\phi\) & Isogenies between elliptic curves. \\
% \(\ker(\varphi)\) & Kernel of the isogeny \(\varphi\). \\
% \(\deg(\varphi)\) & Degree of the isogeny \(\varphi\). \\
% \(\widehat{\varphi}\) & Dual isogeny of \(\varphi\). \\
% \(\End(E)\) & Endomorphism ring of \(E\). \\
% \(\Hom(E,E')\) & Group of isogenies from \(E\) to \(E'\), together with the zero map. \\
% \(j(E)\) & \(j\)-invariant of \(E\). \\
% \(e_m(P,Q)\) & Weil pairing on \(m\)-torsion points. \\
% \(P_0,Q_0\) & Fixed torsion basis on the reference curve \(E_0\). \\
% \(P_\star,Q_\star\) & Torsion basis on the curve \(E_\star\), where \(\star\in\{\mathsf{pk},\mathsf{com},\mathsf{chl},\mathsf{aux}\}\). \\
% \(E_{\mathsf{pk}}\) & Public-key curve. \\
% \(E_{\mathsf{com}}\) & Commitment curve. \\
% \(E_{\mathsf{chl}}\) & Challenge curve. \\
% \(E_{\mathsf{aux}}\) & Auxiliary curve appearing in the response computation. \\
% \(\phi_{\mathsf{sk}}\) & Secret-key isogeny \(E_0\to E_{\mathsf{pk}}\). \\
% \(\phi_{\mathsf{com}}\) & Commitment isogeny \(E_0\to E_{\mathsf{com}}\). \\
% \(\phi_{\mathsf{chl}}\) & Challenge isogeny \(E_{\mathsf{pk}}\to E_{\mathsf{chl}}\). \\
% \(\phi_{\mathsf{rsp}}\) & Response isogeny \(E_{\mathsf{com}}\to E_{\mathsf{chl}}\). \\

\addlinespace[3pt]
\multicolumn{2}{@{}l}{\textbf{Quaternion algebra, orders, and ideals}}\\
\midrule
% \(\Bp\) & Quaternion algebra over \(\Q\) ramified exactly at \(p\) and at infinity. \\
% \(i,j,k\) & Standard generators of \(\Bp\), with
% \(i^2=-1\), \(j^2=-p\), \(ij=-ji\), and \(k=ij\). \\
% % \(\alpha=a+bi+cj+dk\) & General element of \(\Bp\), with \(a,b,c,d\in\Q\). \\
% \(\overline{\alpha}\) & Quaternion conjugate of \(\alpha\). \\
% \(\trd(\alpha)\) & Reduced trace of \(\alpha\). \\
% \(\nrd(\alpha)\) & Reduced norm of \(\alpha\). \\
% \(\calO\) & A quaternion order in \(\Bp\). \\
\(\calO_0\) & The maximal order $\calO_0=\Z\Big\langle 1,\ i,\ \frac{i+j}{2},\ \frac{1+k}{2} \Big\rangle.$ in $\Bp$
\\
% \(\calO_L(I)\) & Left order of \(I\), namely
% \(\{\alpha\in\Bp:\alpha I\subseteq I\}\). \\
% \(\calO_R(I)\) & Right order of \(I\), namely
% \(\{\alpha\in\Bp:I\alpha\subseteq I\}\). \\
% \(\overline{I}\) & Conjugate ideal \(\{\overline{\alpha}:\alpha\in I\}\). \\
% \(\nrd(I)\) & Reduced norm of the ideal \(I\). \\
% \(I\sim J\) & Ideal equivalence; \(I\sim J\) if \(J=I\beta\) for some \(\beta\in\Bp^\times\). \\
% \(I^{-1}\) & Inverse ideal.\\
% ; for an integral invertible ideal, \(I^{-1}=\overline{I}/\nrd(I)\). \\
% \(I_\varphi\) & Connecting ideal associated with a separable isogeny \(\varphi:E\to E'\). \\
% \(E_0[I]\) & Kernel subgroup associated with a left \(\calO_0\)-ideal \(I\). \\
% \(\rho_0\) & Embedding
% \(\Bp\to \End(E_0)\otimes_{\Z}\Q\) sending \(i\mapsto\iota\), \(j\mapsto\pi\). \\
\(\iota,\pi\) & Endomorphisms satisfying \(\iota^2=[-1]\), and \(\pi^2=[-p]\). \\
\(\alpha(P)\) & Action of the endomorphism corresponding to \(\alpha\) on the point \(P\).\\
% \(q_I(\alpha)\) & \[\frac{\nrd(\alpha)}{\nrd(I)}.\]
\(r_\alpha\) & $\min\{r\in\Z_{>0}:r\alpha\in\Z+\Z i+\Z j+\Z k\}$.
\\
\(\chi_I(\beta)\) & An ideal $\frac{\bar{\beta}}{\nrd(I)}I$ equivalent to $I$.
\\

\addlinespace[3pt]
\multicolumn{2}{@{}l}{\textbf{Lattices and reduction}}\\
\midrule
% \(\Lambda\) & Full-rank Euclidean lattice. In the quaternion setting,
% \(\Lambda=\Z\alpha_1+\cdots+\Z\alpha_4\subset\Bp\). \\
% \(\langle x,y\rangle\) & Positive definite bilinear form.
% For quaternions,
% \[
% \langle \alpha,\beta\rangle=\trd(\alpha\overline{\beta}).
% \]
% \\
% \(\|x\|\) & Euclidean norm induced by the bilinear form. For quaternions,
% \(\|\alpha\|^2=2\nrd(\alpha)\). \\
% \(\|\cdot\|_{\nrd}\) & Norm induced by reduced norm, i.e.,
% \(\|\alpha\|_{\nrd}^2=\nrd(\alpha)\). \\
\(\|\cdot\|_\infty\) & Maximum absolute coordinate norm. \\
% \(G_\Lambda\) & Gram matrix of \(\Lambda\). \\
% \(\det(\Lambda)\) & Lattice determinant, equivalently covolume. \\
% \(\det(I)\) & Determinant of the ideal lattice \(I\). \\
% \(\operatorname{covol}_{\nrd}(I)\) & Covolume of \(I\) with respect to the reduced-norm Euclidean structure. \\
% \(\Lambda^\ast\) & Dual lattice of \(\Lambda\) with respect to \(\langle\cdot,\cdot\rangle\). \\
% \(I^\#\) & Trace-dual lattice of a quaternion ideal $I$
% \[
% I^\#=\{\beta\in\Bp:\trd(\alpha\overline{\beta})\in\Z
% \text{ for all } \alpha\in I\}.
% \]
% \\
% \(B_R\) & Euclidean ball of radius \(R\). \\
% \(\lambda_i(\Lambda)\) & \(i\)-th successive minimum of \(\Lambda\). \\
% \(\mu(\Lambda)\) & Covering radius of \(\Lambda\). \\
% \(b_i^\star\) & \(i\)-th Gram--Schmidt vector of an ordered lattice basis. \\
% \(\mu_{i,j}\) & Gram--Schmidt coefficient. \\
% \(\delta\) & LLL reduction parameter; the standard value used here is \(\delta=3/4\). \\
% \(\textsf{LLL}\) & LLL basis-reduction algorithm. \\
% \(\textsf{MLLL}\) & Modified LLL algorithm, which takes a possibly redundant generating set and outputs an LLL-reduced basis of the generated lattice. \\
% \(\textsf{HNF}\) & Hermite normal form, used in previous ideal-normalization routines. \\
\(\lfloor x\rceil\) & Nearest integer to \(x\). \\
$\diamond(x)$ & Nearest fixed-point value to $x$. \\

% \addlinespace[3pt]
% \multicolumn{2}{@{}l}{\textbf{Coordinate matrices and dual computations}}\\
% \midrule
% \(M_I\) & Coordinate matrix of an ideal basis with respect to the ambient basis \((1,i,j,k)\). \\
% \(A_I/r_I\) & Rational coordinate matrix of \(I\), where \(A_I\in\Z^{4\times 4}\) and \(r_I>0\). \\
% \(M/r\) & General numerator-denominator representation of a rational coordinate matrix. \\
% \(r_I\) & Common denominator of the coordinate matrix of \(I\). \\
% \(r_k\) & Denominator-clearing factor for \(I_k\), typically
% \(r_k=\operatorname{lcm}_i r_{\alpha_i}\). \\
% \(M^\#, r^\#\) & Numerator and denominator output by \textsf{LatticeDual}; columns of \(M^\#/r^\#\) generate the dual lattice. \\
% \(\adj(M)\) & Adjugate matrix of \(M\). \\
% \(\Delta\) & Determinant of \(M\), usually \(\Delta=\det(M)\). \\
% \(G_0\) & Gram matrix of \(\nrd\) in the ambient basis \((1,i,j,k)\). \\
% \(\lambda_0\) & Minimum eigenvalue of \(G_0\). \\
% \(C_0\) & Constant in the adjugate bound,
% \(C_0=8\sqrt{\det(G_0)}/\lambda_0^{3/2}=8p\). \\
% \(D_k,s_k\) & Numerator and denominator of the dual basis computed from \(I_k\). \\
% \(s\) & Common denominator, typically \(s=\operatorname{lcm}(s_1,s_2)\). \\
% \(C\) & Concatenated matrix used to generate \(I_1^\#+I_2^\#\). \\
% \(H\) & Matrix output by \textsf{MLLL} from \(C\), so that \(H/s\) generates \(I_1^\#+I_2^\#\). \\
% \(I_\Sigma\) & Sum of trace-dual lattices:
% \[
% I_\Sigma=I_1^\#+I_2^\#.
% \]
% \\
% \(B_\cap\) & Explicit integer-size bound for the compact lattice-intersection algorithm. \\

% \addlinespace[3pt]
% \multicolumn{2}{@{}l}{\textbf{SQIsign keys, signatures, and protocol notation}}\\
% \midrule
% \(\mathsf{pp}\) & Public parameters of SQIsign. \\
% \(\mathsf{pk},\mathsf{sk}\) & Public key and secret key. \\
% % \(\mathsf{msg}\) & Message to be signed. \\
% % \(\sigma\) & Signature. \\
% % \(c\) & Fiat--Shamir challenge. \\
% \(\textsf{HASH}\) & Hash function. \\
% \(\mathsf{hint}_{\mathsf{pk}}\) & Public-key torsion-basis hint. \\
% \(\mathsf{hint}_{\mathsf{aux}},\mathsf{hint}_{\mathsf{chl}}\) & Auxiliary and challenge hints in the signature. \\
% % \(M_{\mathsf{sk}}\) & Secret change-of-basis matrix used in signing. \\
% % \(M_{\mathsf{chl}}\) & Challenge-side change-of-basis matrix included in the signature. \\
% % \(M_\alpha\) & Matrix describing the action of a quaternion element \(\alpha\) on the relevant torsion basis. \\
% % \(I_{\mathsf{sk}}\) & Secret ideal corresponding to the secret-key isogeny. \\
% % \(I_{\mathsf{com}}\) & Commitment ideal. \\
% % \(I'_{\mathsf{chl}}\) & Ideal obtained from the decomposed challenge kernel. \\
% % \(I_{\mathsf{chl}}\) & Challenge ideal transported through the secret ideal. \\
% % \(I_{\mathsf{aux}}\) & Auxiliary ideal used in the response computation. \\
% % \(I_{\mathsf{com,rsp}}\) & Ideal corresponding to the composed commitment-response path. \\
% % \([I_{\mathsf{sk}}]^\ast I'_{\mathsf{chl}}\) & Pullback of the challenge ideal through the secret ideal. \\

\addlinespace[3pt]
\multicolumn{2}{@{}l}{\textbf{Sampling and algorithmic parameters}}\\
\midrule
\(D_{\mathsf{mix}}\) & The smallest prime larger than $2^{4\lambda}$. \\
\(D_{\mathsf{rsp}}\) & SQIsign parameter $2^{\lceil\log_2(\sqrt{p})\rceil}$. \\
\(\mathtt{QUAT\_equiv\_bound\_coeff}\) & \qquad SQIsign parameter $64$. \\
\(\nu_2(x)\) & Dyadic valuation of \(x\). \\
% \(u,v\) & Positive integers used in \textsf{SuitableIdeals}, satisfying a Bezout-type relation. \\
% \(d_1,d_2\) & Normalized norms of the two sampled elements in \textsf{SuitableIdeals}. \\
% \(e,f\) & Two-power exponents appearing in \textsf{SuitableIdeals} and \textsf{IdealToIsogeny}. \\
% \(s,t\) & Indices of the auxiliary parameter ideals \(J_s,J_t\). \\
% \(J_s,J_t\) & SQIsign parameter ideals used inside \textsf{IdealToIsogeny}. \\
% \(B_j\) & Sampling bound for the \(j\)-th coefficient in \textsf{SuitableIdeals}. \\
% \(L\) & List of sampled coefficient vectors in \textsf{SuitableIdeals}. \\

% \addlinespace[3pt]
% \multicolumn{2}{@{}l}{\textbf{Experimental quantities}}\\
% \midrule
% \(R(I)\) & Basis-norm ratio used in the ideal-multiplication experiment:
% \[
% R(I)=\frac{\max_i \nrd(\alpha_i)}{\nrd(I)}.
% \]
% \\
% \(\alpha_1^\#\) & Shortest nonzero vector of the trace-dual lattice \(I^\#\). \\
% \((\alpha_1^\#,\ldots,\alpha_4^\#)\) & Minkowski-reduced basis of \(I^\#\). \\
% \(Y(I)\) & Normalized dual-shortest-vector quantity:
% \[
% Y(I)=\log_2\bigl(N\cdot\nrd(\alpha_1^\#)\bigr).
% \]
% \\
% \(Y_{\min}\) & Minimum of \(Y(I)\) over the experimental samples. \\
\end{longtable}
}

% We denote $r_\alpha:=\min\{r\in\Z : r\alpha\in\Z+\Z i+\Z j+\Z k\}$ for any $\alpha\in\Bp$.

\subsection{Elliptic curves and isogenies}
\label{subsec:isogeny}

We recall the elliptic-curve notions needed to relate the quaternion algorithms used in SQIsign to their geometric meaning.
For more details, see~\cite{Sil09,De17}.
Throughout this subsection, let \(K=\mathbb{F}_{p^2}\), where \(p>3\), and let \(\overline{K}\) denote an algebraic closure of \(K\).
An elliptic curve over \(K\) is a smooth projective curve of genus one with a distinguished point \(0_E\).
When a short Weierstrass model is used, we write
\[
    E : y^2 = x^3 + ax + b,
    \qquad 4a^3+27b^2 \neq 0.
\]
The points of \(E\) form an abelian group with identity \(0_E\).
For an integer \(n\geq 1\), the \(n\)-torsion subgroup is
\[
    E[n] := \{P\in E(\overline{K}) : [n]P = 0_E\}.
\]
If \(\gcd(n,p)=1\), then \(E[n]\cong (\mathbb{Z}/n\mathbb{Z})^2\).
% In SQIsign, torsion bases are used to encode kernels and images of the isogenies appearing in key generation, commitment, challenge, and response computations.

An isogeny \(\phi:E\rightarrow E'\) is a non-constant morphism of elliptic
curves satisfying \(\phi(0_E)=0_{E'}\); equivalently, it is a group homomorphism
with finite kernel. Its degree is denoted by \(\deg(\phi)\), and degrees are
multiplicative under composition. If \(\phi\) is separable, then
\[
    \deg(\phi)=\#\ker(\phi).
\]
Conversely, for every finite subgroup \(G\subset E(\overline{K})\), there is a
quotient curve \(E/G\) and a separable isogeny
\[
    \phi_G:E\longrightarrow E/G
\]
with kernel \(G\). Hence a separable isogeny can be specified by its kernel,
and explicit equations for \(\phi_G\) can be obtained using V\'elu's formulas~\cite{Velu71}. The dual isogeny
\(\widehat{\phi}:E'\rightarrow E\) satisfies
\[
    \widehat{\phi}\circ \phi = [\deg(\phi)]
    \quad\text{and}\quad
    \phi\circ \widehat{\phi} = [\deg(\phi)].
\]

The endomorphism ring of \(E\), denoted \(\operatorname{End}(E)\), is the ring
of isogenies from \(E\) to itself together with the zero map, where addition is
defined pointwise and multiplication is composition. An elliptic curve in
characteristic \(p\) is supersingular if
\[
    \operatorname{End}(E)\otimes_{\mathbb{Z}}\mathbb{Q}
    \cong B_{p,\infty},
\]
where \(B_{p,\infty}\) is the quaternion algebra over \(\mathbb{Q}\) ramified
exactly at \(p\) and at infinity.
% In the SQIsign setting, one fixes a
% supersingular reference curve \(E_0/K\) together with a known maximal order
% \[
%     O_0 \cong \operatorname{End}(E_0)
% \]
% inside \(B_{p,\infty}\). This identification is what allows geometric
% operations on supersingular elliptic curves to be represented by arithmetic
% operations on quaternion orders and ideals.

% We use the following form of the Deuring correspondence. Let \(I\) be an
% integral left \(O_0\)-ideal of reduced norm \(N=\operatorname{nrd}(I)\), with
% \(\gcd(N,p)=1\). It defines a finite subgroup
% \[
%     E_0[I]
%     =
%     \{P\in E_0[N] : \alpha(P)=0_E \text{ for all } \alpha\in I\},
% \]
% and therefore a separable isogeny
% \[
%     \phi_I : E_0 \longrightarrow E_I := E_0/E_0[I]
% \]
% of degree \(N\). The right order
% \[
%     O_R(I)=\{\beta\in B_{p,\infty}: I\beta\subseteq I\}
% \]
% corresponds to the endomorphism ring of the codomain curve \(E_I\). Equivalent
% left ideals define isomorphic codomain curves, while multiplication of
% composable ideals corresponds to composition of the associated isogenies, up to
% the chosen left/right convention \cite{Kohel1996,DLLW2023}.

% This dictionary explains why the algorithms analyzed in this paper can be
% studied on the quaternion side. The routine
% \textsf{RandomIdealGivenNorm} samples ideals representing isogenies of prescribed
% degree, \textsf{RandomEquivalentPrimeIdeal} replaces an ideal by an equivalent
% one with smaller prime norm, and \textsf{IdealToIsogeny} translates the resulting
% ideal data back into curve and torsion-point data. Consequently, bounding the
% integer coefficients that occur in quaternion ideal arithmetic gives a
% fixed-precision bound for the key-generation and signing isogenies used in
% SQIsign \cite{SQISign2020,DLLW2023}.

\subsection{Lattices}
\label{subsec:lattice}

We recall only the lattice notions used in our setting.  Let \(V\) be an
\(n\)-dimensional real vector space equipped with a positive definite
symmetric bilinear form \(\langle \cdot,\cdot\rangle\), and write
\(\|x\|=\sqrt{\langle x,x\rangle}\).  A full-rank lattice in \(V\) is a
discrete subgroup
\[
    \Lambda=\mathbb Z b_1+\cdots+\mathbb Z b_n
\]
generated by an \(\mathbb R\)-basis \(b_1,\ldots,b_n\) of \(V\).  Its Gram
matrix is
\[
    G_\Lambda:=(\langle b_i,b_j\rangle)_{1\le i,j\le n},
\]
and its covolume, or determinant, is
\[
    \det(\Lambda):=\sqrt{\det(G_\Lambda)}.
\]
Equivalently, if \(x=(x_1,\ldots,x_n)^T\in \mathbb Z^n\), then the squared
length of the lattice vector \(\sum_i x_i b_i\) is \(x^TG_\Lambda x\).

The dual lattice of \(\Lambda\), with respect to the same bilinear form, is
\[
    \Lambda^\ast
    =
    \{\, y\in V : \langle x,y\rangle\in \mathbb Z
        \text{ for all } x\in \Lambda \,\}.
:\]
If \(b_1^\ast,\ldots,b_n^\ast\) denotes the basis dual to
\(b_1,\ldots,b_n\), i.e. \(\langle b_i,b_j^\ast\rangle=\delta_{ij}\), then
\(\Lambda^\ast=\mathbb Z b_1^\ast+\cdots+\mathbb Z b_n^\ast\) and
\(\det(\Lambda^\ast)=\det(\Lambda)^{-1}\).

For \(r>0\), let $B_r=\{x\in V:\|x\|\le r\}$.
The \(i\)-th successive minimum of \(\Lambda\) is
\[
    \lambda_i(\Lambda)
    :=
    \inf\{\, r>0 :
    \dim_{\mathbb R}\operatorname{span}_{\mathbb R}(\Lambda\cap B_r)
    \ge i \,\},
    \qquad 1\le i\le n.
\]
Thus \(\lambda_1(\Lambda)\) is the length of a shortest nonzero vector of
\(\Lambda\), and
\[
    0<\lambda_1(\Lambda)\le \cdots \le \lambda_n(\Lambda).
\]
% The covering radius of \(\Lambda\) is
% \[
%     \mu(\Lambda)
%     =
%     \sup_{t\in V}\operatorname{dist}(t,\Lambda),
%     \qquad
%     \operatorname{dist}(t,\Lambda)=\inf_{x\in\Lambda}\|t-x\|.
% \]
% In algorithmic terms, short-vector enumeration with bound \(R\) asks for the
% set \(\Lambda\cap B_R\), while a closest vector to a target \(t\in V\) is an
% element \(x\in\Lambda\) satisfying
% \[
%     \|t-x\|=\operatorname{dist}(t,\Lambda).
% \]

\begin{lemma}[Transference theorem~{\cite[Theorem~2.1]{Ban93}}]
\label{lem:transference}
Let \(\Lambda\subset V\) be a full-rank lattice of rank \(n\).  Then, for
every \(1\le i\le n\),
\[
    1
    \le
    \lambda_i(\Lambda)\lambda_{n-i+1}(\Lambda^\ast)
    \le
    n.
\]
Moreover,
\[
    \frac12
    \le
    \mu(\Lambda)\lambda_1(\Lambda^\ast)
    \le
    \frac n2.
\]
\end{lemma}

The theorem relates the geometry of a lattice to that of its dual: a lattice
cannot have many short independent vectors while its dual also has many short
independent vectors in complementary directions.  In the rank-four lattices
arising from quaternion ideals, the first bound specializes to
\[
    1
    \le
    \lambda_i(\Lambda)\lambda_{5-i}(\Lambda^\ast)
    \le
    4,
    \qquad 1\le i\le 4.
\]

\paragraph{Basis reduction}

We also use LLL-reduced bases~\cite{LLL82}.
Let \(B=(b_1,\ldots,b_n)\) be an ordered basis of a lattice \(\Lambda\), and let
\(b_1^\star,\ldots,b_n^\star\) and \(\mu_{i,j}\) be its Gram--Schmidt
orthogonalization, so that
\[
    b_i=b_i^\star+\sum_{j<i}\mu_{i,j}b_j^\star .
\]
The basis \(B\) is LLL-reduced, with the standard parameter
\(\delta=3/4\), if it is size-reduced,
\[
    |\mu_{i,j}|\le \frac12
    \qquad (1\le j<i\le n),
\]
and satisfies the Lovász condition
\[
    \|b_i^\star+\mu_{i,i-1}b_{i-1}^\star\|^2
    \ge
    \frac34\|b_{i-1}^\star\|^2
    \qquad (2\le i\le n).
\]
% Equivalently,
% \[
%     \|b_i^\star\|^2
%     \ge
%     \Bigl(\frac34-\mu_{i,i-1}^2\Bigr)
%     \|b_{i-1}^\star\|^2 .
% \]
Such a basis is short and nearly orthogonal in a quantitative sense:
\begin{lemma}[LLL basis and successive minima~{\cite[Proposition~1.6]{LLL82}}]
\label{lem:lll-successive-minima}
Let \(\Lambda\) be a rank-\(n\) Euclidean lattice, and let
\(\lambda_i(\Lambda)\) denote its \(i\)-th successive minimum with respect
to the Euclidean norm.  Let
\[
    B=(b_1,\ldots,b_n)
\]
be an LLL-reduced basis of \(\Lambda\). Then, for every \(1\le i\le n\),
\[
    2^{\frac{1-i}{2}}
    \lambda_i(\Lambda)
    \le
    \|b_i\|
    \le
    2^{\frac{n-1}{2}}
    \lambda_i(\Lambda).
\]
% In particular,
% \[
%     \|b_1\|
%     \le
%     2^{\frac{n-1}{2}}\lambda_1(\Lambda)
%     =
%     (\sqrt{2})^{\,n-1}\lambda_1(\Lambda).
% \]
\end{lemma}
% \begin{lemma}[LLL basis and successive minima~{\cite[Proposition~1.6]{LLL82}}]
% \label{lem:lll-successive-minima}
% Let \(\Lambda\) be a rank-\(n\) Euclidean lattice, and let
% \(\lambda_i(\Lambda)\) denote its \(i\)-th successive minimum with respect
% to the Euclidean norm.  Let
% \[
%     B=(b_1,\ldots,b_n)
% \]
% be a \(\delta\)-LLL-reduced basis of \(\Lambda\), where
% \(1/4<\delta\le 1\).
% % Put
% % \[
% %     \alpha_\delta
% %     :=
% %     \frac{1}{\delta-\frac14}
% %     =
% %     \frac{4}{4\delta-1}.
% % \]
% Then, for every \(1\le i\le n\),
% % \[
% %     \alpha_\delta^{(1-i)/2}\lambda_i(\Lambda)
% %     \le
% %     \|b_i\|
% %     \le
% %     \alpha_\delta^{(n-1)/2}\lambda_i(\Lambda).
% % \]
% % Equivalently,
% \[
%     \left(\delta-\frac14\right)^{\frac{i-1}{2}}
%     \lambda_i(\Lambda)
%     \le
%     \|b_i\|
%     \le
%     \left(\frac{1}{\delta-\frac14}\right)^{\frac{n-1}{2}}
%     \lambda_i(\Lambda).
% \]
% In particular,
% \[
%     \|b_1\|
%     \le
%     \alpha_\delta^{(n-1)/2}\lambda_1(\Lambda)
%     =
%     \left(\frac{2}{\sqrt{4\delta-1}}\right)^{n-1}
%     \lambda_1(\Lambda).
% \]
% \end{lemma}


\subsection{Quaternion algebras and Deuring correspondence}
\label{subsec:quaternion}

We recall the quaternion algebras used to represent supersingular isogenies
in SQIsign.
For more details, see~\cite{Voi21, Ler22}.

% We use the quaternion algebra
% \[
%     B_{p,\infty}=\mathbb Q+\mathbb Q i+\mathbb Q j+\mathbb Q ij,
%     \qquad
%     i^2=-1,\quad j^2=-p,\quad ij=-ji,
% \]
% and we also write \(k=ij\). For
% \(\alpha=a+bi+cj+dk\in \Bp\), quaternion conjugation is
% \[
%     \overline{\alpha}=a-bi-cj-dk,
% \]
% and the reduced trace and reduced norm are
% \[
%     \operatorname{trd}(\alpha)=\alpha+\overline{\alpha},
%     \qquad
%     \operatorname{nrd}(\alpha)=\alpha\overline{\alpha}.
% \]
% An order \(\mathcal O\subset \Bp\) is a rank-four \(\mathbb Z\)-lattice that is
% also a subring containing \(1\), and it is maximal if it is not properly contained
% in any larger order. A left \(\mathcal O\)-ideal is a finitely generated
% \(\mathbb Z\)-lattice \(I\subset \Bp\) satisfying \(\mathcal O I\subseteq I\).
% Its right order is
% \[
%     O_R(I)=\{\beta\in \Bp : I\beta\subseteq I\}.
% \]
% For an invertible integral left \(\mathcal O\)-ideal \(I\), the reduced norm
% \(\operatorname{nrd}(I)\) is the positive integer satisfying
% \[
%     I\overline{I}=\operatorname{nrd}(I)\mathcal O.
% \]
% Two left ideals \(I\) and \(J\) are equivalent if \(J=I\gamma\) for some
% \(\gamma\in \Bp^\times\).

% \paragraph{Deuring correspondence.}
% For a supersingular elliptic curve \(E\) defined over \(\Fpp\), the
% endomorphism ring \(\operatorname{End}(E)\) is isomorphic to a maximal order
% in \(B_{p,\infty}\)~\cite{Deu41}. Conversely, every maximal order in
% \(B_{p,\infty}\) occurs as the endomorphism ring of some supersingular
% elliptic curve. Thus, on objects, the Deuring correspondence identifies
% supersingular elliptic curves with maximal orders in \(B_{p,\infty}\), up to
% the usual isomorphism or type-class equivalence.

% The correspondence also relates isogenies to connecting ideals. Let
% \[
%     \varphi:E\longrightarrow E'
% \]
% be a separable isogeny, and set
% \[
%     \calO=\operatorname{End}(E),
%     \qquad
%     \calO'=\operatorname{End}(E').
% \]
% After identifying \(\calO\) and \(\calO'\) with maximal orders in
% \(B_{p,\infty}\), the isogeny \(\varphi\) determines a connecting ideal
% \[
%     I_\varphi
%     :=
%     \{\alpha\circ\varphi:\alpha\in \operatorname{Hom}(E',E)\}.
% \]
% This ideal has reduced norm equal to the degree of the isogeny:
% \[
%     \operatorname{nrd}(I_\varphi)=\deg(\varphi).
% \]
% Equivalent isogenies give equivalent ideals, the dual isogeny corresponds to
% the conjugate ideal,
% \[
%     I_{\widehat{\varphi}}=\overline{I_\varphi},
% \]
% and, with the convention used here, for composable isogenies
% \(\varphi:E\to E'\) and \(\psi:E'\to E''\), one has
% \[
%     I_{\psi\circ\varphi}=I_\varphi I_\psi.
% \]

% For the reference curve used in SQIsign, let
% \[
%     E_0: y^2=x^3+x
% \]
% over \(\Fpp\), with \(p\equiv 3 \pmod 4\). Choose
% \(\sqrt{-1}\in\mathbb F_{p^2}\) such that \((\sqrt{-1})^2=-1\), and define
% endomorphisms of \(E_0\) by
% \[
%     \iota(x,y):=(-x,\sqrt{-1}y),
%     \qquad
%     \pi(x,y):=(x^p,y^p).
% \]
% Then
% \[
%     \iota^2=[-1],
%     \qquad
%     \pi^2=[-p],
%     \qquad
%     \iota\pi=-\pi\iota.
% \]
% Hence the map
% \[
%     \rho_0:
%     B_{p,\infty}
%     \longrightarrow
%     \operatorname{End}(E_0)\otimes_{\mathbb Z}\mathbb Q,
%     \qquad
%     i\mapsto \iota,\quad
%     j\mapsto \pi,\quad
%     k=ij\mapsto \iota\pi,
% \]
% restricts to an isomorphism~\cite{Ibu82}
% \[
%     \rho_0|_{\calO_0}:
%     \calO_0:=
%     \mathbb Z\Big\langle
%         1,\,
%         i,\,
%         \frac{i+j}{2},\,
%         \frac{1+k}{2}
%     \Big\rangle
%     \xrightarrow{\sim}
%     \operatorname{End}(E_0).
% \]
% Thus a left \(\calO_0\)-ideal \(I\) of norm \(N\), with \(\gcd(N,p)=1\),
% defines the subgroup
% \[
%     E_0[I]
%     =
%     \{P\in E_0[N]:
%       \alpha(P)=0 \text{ for all } \alpha\in I\},
% \]
% and therefore a separable isogeny
% \[
%     \phi_I:E_0\longrightarrow E_0/E_0[I]
% \]
% of degree \(N=\operatorname{nrd}(I)\). Moreover, the endomorphism ring of the
% codomain is isomorphic to the right order \(O_R(I)\). This is the direction of
% the correspondence used by the ideal-to-isogeny translation routines in
% SQIsign.

% The following table summarizes the dictionary used throughout the paper.

% \begin{center}
% \begin{tabular}{@{}ll@{}}
% \toprule
% \textbf{Elliptic-curve world} & \textbf{Quaternion-algebra world} \\
% \midrule
% Supersingular elliptic curve \(E\) &
% \(\calO=\operatorname{End}(E)\subset B_{p,\infty}\) \\[2pt]
% Separable isogeny \(\varphi:E\to E'\) &
% \(I_\varphi\), a \(\calO,\calO'\)-connecting ideal \\[2pt]
% \(\deg(\varphi)\) &
% \(\operatorname{nrd}(I_\varphi)\) \\[2pt]
% Equivalent isogenies &
% Equivalent ideals \(I_\varphi\sim I_\psi\) \\[2pt]
% Dual isogeny \(\widehat{\varphi}:E'\to E\) &
% Conjugate ideal \(I_{\widehat{\varphi}}=\overline{I_\varphi}\) \\[2pt]
% Composition \(\psi\circ\varphi\) &
% Ideal product \(I_\varphi I_\psi\) \\
% \bottomrule
% \end{tabular}
% \end{center}

\paragraph{Quaternion algebras}

Let $B_{p, \infty}$ denote a quaternion algebra over $\Q$
% ramified at $p$ and $\infty$.
% That is, at a place $v$ of $\Q$, $B_{p,\infty} \otimes_\Q \Q_v$ is not isomorphic to $M_{2\times 2}(\Q)$.
whose $\Q$-basis is $\{1, i, j, k:=ij\}$ with $i^2=-1, j^2=-p, ij=-ji$;
% We use the representation $B_{p,\infty} \cong \Q\langle{i,j}\rangle$ where $i^2 = -1$ and $j^2 = -p$.
every quaternion element can be represented by $\alpha = a + b i +cj+dk \in B_{p,\infty}$ for some $a,b,c,d \in \Q$.
The canonical conjugate $\bar{\alpha}$ is defined as $a- bi-cj-dk$.
The reduced trace and reduced norm are defined by
$$
    \mathrm{trd}(\alpha):=\alpha + \bar{\alpha}, \; \mathrm{nrd}(\alpha):=\alpha \bar{\alpha} = \bar{\alpha}\alpha \in \Q.
$$
A \textit{fractional ideal} in $B_{p,\infty}$ is a $\Z$-submodule of rank 4.
An \textit{order} $\calO$ is a fractional ideal that is also a subring of $\Bp$.
An order $\calO$ is \textit{maximal} if it is not properly contained in any other order. A fractional ideal $I$ is called a left $\calO$-ideal if $\alpha I \subseteq I$ for all $\alpha \in \calO$. It is called a right $\calO$-ideal if $I \alpha \subseteq I$ for all $\alpha \in \calO$.
For a fractional ideal \(I\), its left and right orders are defined by
\[
O_L(I):=\{\alpha\in B_{p,\infty}:\alpha I\subseteq I\},
\qquad
O_R(I):=\{\alpha\in B_{p,\infty}:I\alpha\subseteq I\}.
\]
Its conjugate is defined by $\bar{I} := \{\bar{\alpha} : \alpha\in I\}$ and its reduced norm is defined by $\nrd(I) := \gcd\{\nrd(\alpha) : \alpha\in I\}$.
Also, its inverse ideal is defined by $I^{-1}:=\{x\in B_{p,\infty}: Ix\subseteq O_L(I)\}$, which satisfies $I^{-1}=\frac{\bar{I}}{\nrd(I)}$ if $O_L(I)$ is maximal (or, equivalently, $O_R(I)$ is).

We say that $I$ is integral if it is contained in its left order (or, equivalently, to its right order).
Moreover, $I$ is said to be cyclic if for all prime $\ell$ we have that $I\not\subset\ell O_L(I)$.
Two left \(\calO\)-ideals \(I\) and \(J\) are said to be \emph{equivalent}, written
\(I\sim J\), if there exists \(\beta\in B_{p,\infty}^{\times}\) such that $J=I\beta$.
In this case,
\[
O_L(J)=O_L(I)
\qquad\text{and}\qquad
O_R(J)=\beta^{-1}O_R(I)\beta.
\]
The corresponding equivalence classes form the \emph{left ideal class set} of \(\calO\),
which is the object that appears on the quaternionic side of the Deuring correspondence.

\paragraph{Deuring correspondence.}

For a supersingular elliptic curve \(E\) defined over $\Fpp$, the endomorphism ring
\(\operatorname{End}(E)\) is isomorphic to a maximal order in \(B_{p,\infty}\)~\cite{Deu41}.
Conversely, every maximal order in \(B_{p,\infty}\) occurs as
\(\operatorname{End}(E)\) for some supersingular elliptic curve \(E\).
Thus, on objects, the Deuring correspondence identifies isomorphism classes of supersingular elliptic curves with isomorphism classes of maximal orders in \(B_{p,\infty}\).

For a separable isogeny
\(\varphi:E\to E'\), we define the associated connecting ideal by
\[
I_\varphi:=\{\alpha\circ\varphi:\alpha\in \operatorname{Hom}(E',E)\}.
\]
% viewed inside \(B_{p,\infty}\) via \(\rho\).
% Then
% % \(I_\varphi\) is a left
% % \(\mathcal O\)-ideal and a right \(\mathcal O'\)-ideal, and
% \[
% \operatorname{nrd}(I_\varphi)=\deg(\varphi).
% \]
% Equivalent isogenies give equivalent ideals, the dual isogeny corresponds to the
% conjugate ideal,
% \[
% I_{\widehat{\varphi}}=\overline{I_\varphi},
% \]
% and for composable isogenies \(\varphi:E\to E'\) and \(\psi:E'\to E''\),
% \[
% I_{\psi\circ\varphi}=I_\varphi I_\psi.
% \]

Choose \(\sqrt{-1}\in\mathbb{F}_{p^2}\) such
that \(\sqrt{-1}^2=-1\), and define endomorphisms in $\End(E_0)$
\[
\iota(x,y):=(-x,\sqrt{-1} y),
\qquad
\pi(x,y):=(x^p,y^p).
\]
Then \(\iota^2=[-1]\), \(\pi^2=[-p]\), and \(\iota\pi=-\pi\iota\). The map
\[
\rho_0:\ B_{p,\infty}\longrightarrow
\operatorname{End}(E_0)\otimes_{\mathbb{Z}}\mathbb{Q},
\qquad
i\mapsto \iota,\quad j\mapsto \pi,\quad k=ij\mapsto \iota\pi,
\]
restricts to an isomorphism~\cite{Ibu82}
\[
\rho_0|_{\calO_0}:\ \mathcal O_0:=
\mathbb{Z}\Big\langle 1,\, i,\, \frac{i+j}{2},\, \frac{1+k}{2}\Big\rangle
\xrightarrow{\sim} \operatorname{End}(E_0).
\]

We give a table which summarizes the Deuring correspondence, the correspondence between objects of elliptic curves and objects of quaternion algebras.

\begin{center}
\begin{tabular}{@{}l@{\hspace{15pt}}l@{}}
\toprule
\textbf{Elliptic-curve world} & \textbf{Quaternion-algebra world} \\
\midrule
Supersingular elliptic curve \(E\) & \(\mathcal O=\operatorname{End}(E)\subset B_{p,\infty}\) \\[2pt]
\(\varphi:E\to E'\) & \(I_\varphi\), a \(\mathcal O\),\(\mathcal O'=\operatorname{End}(E')\)-connecting ideal \\[2pt]
\(\deg(\varphi)\) & \(\operatorname{nrd}(I_\varphi)\) \\[2pt]
\(\varphi,\psi:E\to E'\) & \(I_\varphi\sim I_\psi\) \\[2pt]
\(\widehat{\varphi}:E'\to E\) & \(I_{\widehat{\varphi}}=\overline{I_\varphi}\) \\[2pt]
\(\psi\circ\varphi\) & \(I_\varphi\cdot I_\psi\) \\
\bottomrule
\end{tabular}
\end{center}

\paragraph{Quaternion lattices.}

For fixed-precision analysis, these ideals are viewed as rank-four Euclidean
lattices. The trace pairing
\[
    \langle \alpha,\beta\rangle
    =
    \operatorname{trd}(\alpha\overline{\beta})
\]
is positive definite on \(B_{p,\infty}\otimes_{\mathbb Q}\mathbb R\), and
\[
    \|\alpha\|^2
    :=
    \langle \alpha,\alpha\rangle
    =
    2\operatorname{nrd}(\alpha).
\]
Hence a quaternion lattice $\Lambda
    = \mathbb Z\alpha_1+\mathbb Z\alpha_2
    +\mathbb Z\alpha_3+\mathbb Z\alpha_4$
with \(\mathbb Q\)-linearly independent \(\alpha_i\in B_{p,\infty}\) is a
rank-four Euclidean lattice with Gram matrix
\[
    G_\Lambda
    =
    \bigl(\operatorname{trd}(\alpha_i\overline{\alpha_j})\bigr)_{1\le i,j\le 4}\enspace.
\]
Thus, for \(x=(x_1,\ldots,x_4)^T\in\mathbb Z^4\),
\[
    \left\|\sum_{i=1}^4 x_i\alpha_i\right\|^2
    =
    x^TG_\Lambda x.
\]
Consequently, enumerating or sampling quaternion elements of bounded reduced
norm is equivalent to enumerating or sampling vectors in this positive
definite lattice with respect to the quadratic form \(x^TG_\Lambda x\).

Throughout this paper, a quaternion ideal is represented by a coordinate
matrix with respect to the fixed ambient basis \((1,i,j,k)\) of
\(B_{p,\infty}\). More precisely, for a basis $I=\mathbb Z\alpha_1+\mathbb Z\alpha_2+\mathbb Z\alpha_3+\mathbb Z\alpha_4$,
we write each basis element as
\[
    \alpha_\ell
    =
    M_{1,\ell}
    +M_{2,\ell}i
    +M_{3,\ell}j
    +M_{4,\ell}k
    \qquad (1\leq \ell\leq 4),
\]
and define its coordinate matrix by
\[
    M_I:=(M_{m,\ell})_{1\leq m,\ell\leq 4}.
\]
Thus the \(\ell\)-th column of \(M_I\) is the coordinate vector of
\(\alpha_\ell\), and
\[
    (\alpha_1,\alpha_2,\alpha_3,\alpha_4)
    =
    (1,i,j,k)M_I.
\]
When the coefficients are rational, we write \(M_I=A_I/r_I\), where
\(A_I\in\mathbb Z^{4\times 4}\) and \(r_I>0\), so that
\[
    (\alpha_1,\alpha_2,\alpha_3,\alpha_4)
    =
    (1,i,j,k)\frac{A_I}{r_I}.
\]

We also define the trace-dual lattice
\[
    I^\#
    :=
    \{\beta\in \Bp :
      \operatorname{trd}(\alpha\overline{\beta})\in\mathbb Z
      \text{ for all } \alpha\in I\},
\]
which is the dual lattice of \(I\) with respect to the trace pairing.

For any maximal order $\calO$ in $\Bp$ and any left $\calO$-ideal $I$, we define \[q_I: I \to \Z,\ \alpha\mapsto\frac{\nrd(\alpha)}{\nrd(I)}.\]

% \begin{lemma}[{\cite[Lemma~48]{DLRW24}}]
% \label{lem:minkowski-norm}
% Let $I$ be a left-ideal of a maximal order $\calO$ in $\Bp$ and $(b_1,\dots,b_4)$ be a Minkowski-reduced basis of $I$. Then,
% \[ q_I(b_4) \le \frac{8p}{\pi^2}. \]
% \end{lemma}
Now we state some lemmas used repeatedly in the later fixed-precision analysis.

\begin{lemma}
\label{lem:lll-norm}
Let $I$ be a left-ideal of a maximal order $\calO$ in $\Bp$, and $(b_1,\dots,b_4)$ be its LLL-reduced basis. Then,
\[ \|b_i\|^2 \le \frac{64p}{\pi^2}\nrd(I). \]
\end{lemma}

\begin{proof}
Let $(b_1^\star,\dots,b_4^\star)$ be a Minkowski-reduced basis of $I$, and $\lambda_i$ be its $i$-th successive minima for each $i$.
Then, as a well-known fact from~\cite{WL56}, we have
\[ \|b_i^\star\| = \lambda_i. \quad(1\le i\le4) \]
By applying this equality and~\cite[Lemma~48]{DLRW24} to Lemma~\ref{lem:lll-successive-minima} with $n=4$, we conclude the statement.\\
\end{proof}

\begin{lemma}[{\cite[Lemma~5]{KLKL25}}]
\label{lem:denombound}
Let $\mathcal{O}$ be a maximal order in $B_{p,\infty}$ with a $\mathcal{O}_0,\mathcal{O}$-connecting ideal $I$ s.t. $\mathrm{nrd}(I) = N$. Then, for any $\alpha \in \mathcal{O}$, $r_{\alpha} \leq 2N$.
\end{lemma}

\begin{lemma}[{\cite[Lemma~9]{KLKL25}}]
\label{lem:normintersection}
Let $\mathcal{O}$ be a maximal order in $B_{p,\infty}$ and $I_1,I_2$ be intersection of two left $\mathcal{O}$-ideals. Then,
\[ \mathrm{nrd}(I_1 \cap I_2) \leq \mathrm{nrd}(I_1)\mathrm{nrd}(I_2). \]
\end{lemma}

\begin{lemma}[{\cite[Lemma~15]{KLKL25}}]
\label{lem:detnrd}
Let $I$ be a left/right-ideal of a maximal order $\mathcal{O}$ in $B_{p,\infty}$. Then,
\[ \det(I) = \det(\mathcal{O})\mathrm{nrd}(I)^2. \]
\end{lemma}

\begin{lemma}[{\cite[Lemma~16]{KLKL25}}]
\label{lem:detO0}
$\det(\mathcal{O}_0) = \frac{p}{4}$
\end{lemma}

% Lemma~\ref{lem:lll-norm} converts reduced
% ideal norms into bounds on reduced bases. Lemma~\ref{lem:denombound} controls
% the denominators of quaternion elements represented in a fixed integral basis.
% Lemma~\ref{lem:normintersection} bounds the norm growth caused by ideal
% intersection, while Lemma~\ref{lem:detnrd} and Lemma~\ref{lem:detO0} translate
% between ideal norms and lattice determinants. These conversions allow the
% integer growth of ideal multiplication, ideal intersection, random equivalent
% ideal sampling, and the ideal-to-isogeny translation to be bounded uniformly.


% \subsection{Quaternion algebras and Deuring correspondence}
% \label{subsec:quaternion}

% In SQIsign, the relevant lattices arise from quaternion orders and ideals.
% Let
% \[
%     B_{p,\infty}=\mathbb Q+\mathbb Q i+\mathbb Q j+\mathbb Q ij,
%     \qquad
%     i^2=-1,\quad j^2=-p,\quad ij=-ji,
% \]
% and let \(\overline{\alpha}\), \(\operatorname{trd}(\alpha)\), and
% \(\operatorname{nrd}(\alpha)\) denote quaternion conjugation, reduced trace,
% and reduced norm, respectively.  The form
% \[
%     \langle \alpha,\beta\rangle
%     =
%     \operatorname{trd}(\alpha\overline{\beta})
% \]
% is positive definite on \(B_{p,\infty}\otimes_{\mathbb Q}\mathbb R\), and
% \[
%     \|\alpha\|^2
%     =
%     \langle \alpha,\alpha\rangle
%     =
%     2\operatorname{nrd}(\alpha).
% \]
% Hence a quaternion lattice
% \[
%     \Lambda
%     =
%     \mathbb Z\alpha_1+\mathbb Z\alpha_2
%     +\mathbb Z\alpha_3+\mathbb Z\alpha_4
% \]
% with \(\mathbb Q\)-linearly independent \(\alpha_i\in B_{p,\infty}\) is a
% rank-four Euclidean lattice with Gram matrix
% \[
%     G_\Lambda
%     =
%     \bigl(\operatorname{trd}(\alpha_i\overline{\alpha_j})\bigr)_{1\le i,j\le 4}.
% \]
% Thus, for \(x=(x_1,\ldots,x_4)^T\in\mathbb Z^4\),
% \[
%     \left\|\sum_{i=1}^4 x_i\alpha_i\right\|^2
%     =
%     x^TG_\Lambda x.
% \]
% Consequently, enumerating or sampling quaternion elements of bounded reduced
% norm is equivalent to enumerating or sampling vectors in this positive
% definite lattice with respect to the quadratic form \(x^TG_\Lambda x\).

% For any maximal order $\calO$ in $\Bp$ and any left $\calO$-ideal $I$, we define \[q_I: I \to \Z,\ \alpha\mapsto\frac{\nrd(\alpha)}{\nrd(I)}.\]

% \begin{lemma}[{\cite[Lemma~48]{DLRW24}}]
% \label{lem:minkowski-norm}
% Let $I$ be a left-ideal of a maximal order $\calO$ in $\Bp$ and $(b_1,\dots,b_4)$ be a Minkowski-reduced basis of $I$. Then,
% \[ q_I(b_4) \le \frac{8p}{\pi^2}. \]
% \end{lemma}

% \begin{lemma}
% \label{lem:lll-norm}
% Let $I$ be a left-ideal of a maximal order $\calO$ in $\Bp$, and $(b_1,\dots,b_4)$ be its $\delta$-LLL reduced basis. Then,
% \[ \|b_i\|^2 \le \left(\frac{1}{\delta-\frac14}\right)^3\frac{8p}{\pi^2}\nrd(I). \]
% \end{lemma}

% \begin{proof}
% Let $(b_1^\star,\dots,b_4^\star)$ be a Minkowski-reduced basis of $I$, and $\lambda_i$ be its $i$-th successive minima for each $i$. Then, as a well-known fact from~\cite{WL56}, we have
% \[ \|b_i^\star\| = \lambda_i. \quad(1\le i\le4) \]
% By applying this equality and Lemma~\ref{lem:minkowski-norm} to the inequality
% \[ \|b_i\|^2 \le \left(\frac{1}{\delta-\frac14}\right)^3\lambda_i, \quad(1\le i\le4) \]
% we conclude the statement.\\
% \end{proof}

% \begin{lemma}[{\cite[Lemma~5]{KLKL25}}]
% \label{lem:denombound}
% Let $\mathcal{O}$ be a maximal order in $B_{p,\infty}$ with a $\mathcal{O}_0,\mathcal{O}$-connecting ideal $I$ s.t. $\mathrm{nrd}(I) = N$. Then, for any $\alpha \in \mathcal{O}$, $r_{\alpha} \leq 2N$.
% \end{lemma}

% \begin{lemma}[{\cite[Lemma~9]{KLKL25}}]
% \label{lem:normintersection}
% Let $\mathcal{O}$ be a maximal order in $B_{p,\infty}$ and $I_1,I_2$ be intersection of two left $\mathcal{O}$-ideals. Then,
% \[ \mathrm{nrd}(I_1 \cap I_2) \leq \mathrm{nrd}(I_1)\mathrm{nrd}(I_2). \]
% \end{lemma}

% \begin{lemma}[{\cite[Lemma~15]{KLKL25}}]
% \label{lem:detnrd}
% Let $I$ be a left/right-ideal of a maximal order $\mathcal{O}$ in $B_{p,\infty}$. Then,
% \[ \det(I) = \det(\mathcal{O})\mathrm{nrd}(I)^2. \]
% \end{lemma}

% \begin{lemma}[{\cite[Lemma~16]{KLKL25}}]
% \label{lem:detO0}
% $\det(\mathcal{O}_0) = \frac{p}{4}$
% \end{lemma}

\subsection{SQIsign}

SQIsign~\cite{DKL+20} is an isogeny-based digital signature scheme over supersingular elliptic curves.
Conceptually, it is obtained by applying the Fiat--Shamir transform to a proof of knowledge of hidden endomorphism information associated with a public supersingular curve. For more details, see~\cite{AAA+25}.
% Throughout, let \(E \to E'\) denote an isogeny between supersingular elliptic curves over \(\mathbb{F}_{p^2}\), let \(j(E)\) denote the \(j\)-invariant of \(E\), and let \(\operatorname{End}(E)\) denote the endomorphism ring of \(E\).

% \algobox{
% \textbf{Hard problem.}
% \begin{enumerate}
%   \setlength{\itemsep}{0.15em}
%   \item The underlying computational assumption is the hardness of the
%   \emph{supersingular endomorphism ring problem}: given a supersingular curve
%   \(E\), compute \(\operatorname{End}(E)\).
%   \item Equivalently at a high level, this is tied to the supersingular
%   \emph{isogeny path problem}: given two supersingular curves \(E_1,E_2\), find
%   an isogeny \(E_1\to E_2\).
%   \item SQIsign hides the endomorphism structure of the public key curve
%   \(E_{\mathsf{pk}}\). A signature is a compact non-interactive proof that the
%   signer can answer a hash-determined isogeny challenge using this hidden
%   structure.
% \end{enumerate}
% }

SQIsign is a digital signature scheme specified by four algorithms
\[
    \mathsf{SQIsign}
    =
    (\mathsf{Setup},\mathsf{KeyGen},\mathsf{Sign},\mathsf{Verify}),
\]
where \(\mathsf{Setup}\) generates public parameters, \(\mathsf{KeyGen}\) outputs a public/secret key pair, \(\mathsf{Sign}\) produces a signature on a message, and \(\mathsf{Verify}\) checks its validity.
At a high-level, each algorithm can be explained as follows:

\begin{algobox}
\underline{\textbf{\(\mathsf{Setup}(1^\lambda)\)}}
\medskip
\begin{enumerate}
  \setlength{\itemsep}{0.15em}
  \item Choose public parameters for the security level \(\lambda\), including
  a prime \(p\).
  % a fixed supersingular reference curve
  % \(E_0/\mathbb{F}_{p^2}\), and a known reference order
  % \(\operatorname{End}(E_0)\).
  \item Fix the prescribed degrees and encodings for key, commitment,
  challenge, and response isogenies, together with torsion-basis conventions.
  \item Fix a hash function \(H\), modeled as a random oracle in the
  Fiat--Shamir view, and output the public parameters \(\mathsf{pp}\).
\end{enumerate}
\end{algobox}

\begin{algobox}
\underline{\textbf{\(\mathsf{KeyGen}(\mathsf{pp})\)}}
\medskip
\begin{enumerate}
  \setlength{\itemsep}{0.15em}
  \item Sample a secret isogeny
  \[
      \varphi_{\mathsf{sk}}:E_0\longrightarrow E_{\mathsf{pk}},
  \]
  internally represented through the quaternion-ideal machinery.
  \item Set the public key to be the codomain curve, together with auxiliary
  public decoding data:
  \[
      \mathsf{pk}=(E_{\mathsf{pk}},\mathsf{hint}_{\mathsf{pk}}).
  \]
  \item Store as secret key a compact representation of
  \(\varphi_{\mathsf{sk}}\), equivalently enough information to recover the
  hidden endomorphism structure of \(E_{\mathsf{pk}}\), together with auxiliary
  basis data used for signing.
\end{enumerate}
\end{algobox}

\begin{algobox}
\underline{\textbf{\(\mathsf{Sign}(\mathsf{sk},m)\).}}
\medskip
\begin{enumerate}
  \setlength{\itemsep}{0.15em}
  \item \emph{Commitment.} Sample a fresh random commitment isogeny
  \[
      \varphi_{\mathsf{com}}:E_0\longrightarrow E_{\mathsf{com}}.
  \]
  \item \emph{Challenge.} Compute the Fiat--Shamir challenge
  \[
      c := H(\mathsf{pk}\,\|\,j(E_{\mathsf{com}})\,\|\,m),
  \]
  and interpret \(c\) as defining a challenge isogeny
  \[
      \varphi_{\mathsf{chl}}:E_{\mathsf{pk}}\longrightarrow E_{\mathsf{chl}}.
  \]
  \item \emph{Response.} Using the secret isogeny
  \(\varphi_{\mathsf{sk}}\) and the commitment data, compute a compact
  descriptor \(R\) of a response isogeny
  \[
      \varphi_{\mathsf{rsp}}:E_{\mathsf{com}}\longrightarrow E_{\mathsf{chl}},
  \]
  subject to SQIsign's no-backtracking condition.
  \item Output
  \[
      \sigma=(c,R).
  \]
  In the concrete scheme, \(R\) is encoded using auxiliary curves, small
  integers, torsion-basis hints, and change-of-basis matrices rather than by
  explicitly transmitting the full isogeny.
\end{enumerate}
\end{algobox}

\begin{algobox}
\underline{\textbf{\(\mathsf{Verify}(\mathsf{pk},m,\sigma)\)}}
\medskip
\begin{enumerate}
  \setlength{\itemsep}{0.15em}
  \item Parse \(\sigma=(c,R)\). Reject if the encoded curves, integers, hints,
  or matrices are malformed, or if the relevant curves are not supersingular.
  \item From \(\mathsf{pk}\) and \(c\), reconstruct the challenge isogeny
  \[
      \varphi_{\mathsf{chl}}:E_{\mathsf{pk}}\longrightarrow E_{\mathsf{chl}}.
  \]
  \item From the response descriptor \(R\), evaluate the claimed response
  structure and recover the corresponding commitment curve
  \(E_{\mathsf{com}}\) up to isomorphism.
  \item Check that the response data is consistent with an isogeny
  \(E_{\mathsf{com}}\to E_{\mathsf{chl}}\) and satisfies the required
  no-backtracking and kernel-validity conditions.
  \item Recompute
  \[
      c' := H(\mathsf{pk}\,\|\,j(E_{\mathsf{com}})\,\|\,m),
  \]
  and accept iff \(c'=c\) and all previous consistency checks pass.
\end{enumerate}
\end{algobox}